\section{Selmer seed for the reducible split-pair anchor}\label{sec:reducible}

We prove Proposition~\ref{prop:anchor} for a non-CM elliptic curve
\(E/\mathbb Q\) with \(E(\mathbb Q)[2]\ne0\).  The construction has two
stages.  First, finite two-descent produces a rational twist, possibly
followed by an isogeny, whose ordinary two-Selmer group consists only of
torsion classes.  A further imaginary twist then gives ordinary Selmer
corank one over a split quadratic field.  The determinant comparison of
Sections~\ref{sec:split-measures} and~\ref{sec:split-local} supplies the
analytic conclusion and the exact leading-coefficient identity only after
these finite constructions are complete.

All dimensions in this section are over \(\mathbb F_2\).  We call
\[
 d_2(C)=\dim\operatorname{Sel}_2(C/\mathbb Q)-\dim C(\mathbb Q)[2]
\]
the \emph{pure dimension}.  Twisting factors will be products of distinct
signed primes
\[
 p^*=(-1)^{(p-1)/2}p,\qquad p\nmid 2N_E.
\]
Such a product is an odd fundamental discriminant; the empty product,
interpreted as \(1\), is also allowed unless explicitly excluded.

\begin{proposition}[Finite Selmer seed]\label{prop:reducible-seed}
There are a permitted twisting factor \(h\) and an elliptic curve \(E_0\),
isogenous over \(\mathbb Q\) to \(E^h\), with rational two-torsion and
\[
                    \dim{\rm Sel}_2(E_0)=\dim E_0(\mathbb Q)[2].
                    \tag{seed}\label{eq:seed}
\]
\end{proposition}

The proposition will follow from the full-rational-plane construction and
the one-line construction below.  Neither construction uses analytic
nonvanishing or finiteness of the Tate--Shafarevich group.

\subsection{Labels and finite changes of local conditions}

At each stage let \(S\) contain \(2,\infty\), all bad primes of the current
curve, the original excluded primes, and all previously used twisting
primes.  We may enlarge \(S\), for example to contain the support of the
coefficients of a chosen equation.  Set
\[
 \begin{split}
 H=H(S)&=\langle-1,q:q\in S\text{ finite}\rangle
               \subset\mathbb Q^\times/\mathbb Q^{\times2},\\
 m&=\dim H=1+\#\{q\in S:q\text{ finite}\}=|S|,\\
 V&=\bigoplus_{v\in S}\mathbb Q_v^\times/\mathbb Q_v^{\times2}.
 \end{split}
\]
A fresh prime \(p_i\) has a \emph{label} \(l_i\in H^*\), given by evaluation
of the radical characters at its Frobenius.  Write \(e_{ij}\) for the
additive Legendre symbol of \(p_j^*\) at \(p_i\), for \(i\ne j\).

\begin{lemma}[Realization of labels and edges]\label{lem:reducible-labels}
There is a linear map \(\xi:H^*\to V\) such that
\(\operatorname{loc}_S(p_i^*)=\xi(l_i)\), and the map
\[
 (z,y)\longmapsto\operatorname{loc}_S(z)+\xi(y),
 \qquad H\oplus H^*\longrightarrow V,
\]
is injective.  Arbitrary finite label lists and edge values are realizable
by distinct fresh primes, subject only to
\[
 e_{ij}+e_{ji}=l_i(-1)l_j(-1).
\]
\end{lemma}

\begin{proof}
Quadratic reciprocity determines the localization of \(p_i^*\) at each
odd place of \(S\).  At \(2\), its squareclass is determined by the symbol
of \(2\), since \(p_i^*\equiv1\pmod4\); its real sign is determined by the
symbol of \(-1\).  These descriptions are linear in \(l_i\), giving
\(\xi\).  Global Hilbert reciprocity gives
\[
 \langle\operatorname{loc}_S(z),\xi(y)\rangle=y(z)
 \qquad(z\in H,\ y\in H^*).
\]
The subspace \(\operatorname{loc}_S H\) is isotropic: outside \(S\), both
arguments are odd local units and their Hilbert symbols are trivial.
If \(\operatorname{loc}_S(z)+\xi(y)=0\), pairing with this subspace first
gives \(y=0\), and valuations and sign then give \(z=0\).
Notice also that \(\dim V=2m\).

The displayed edge relation is quadratic reciprocity in signed-prime
notation.  Choose the primes successively.  Their labels prescribe
congruence classes at the fixed support, while the edges to earlier primes
prescribe additional quadratic-residue classes.  The Chinese remainder
theorem and Dirichlet's theorem realize these prescriptions; reciprocity
then determines precisely the reverse edges.  This also realizes every
label list.
\end{proof}

We identify two-torsion modules across quadratic twists.  Local duality
makes their Kummer conditions self-orthogonal for the Weil pairing;
conditions for dual isogenies are orthogonal complements.  At infinity we
use the usual \(H^1\).  At a good odd prime not ramified in a twist, the
Kummer condition is the full unramified condition.

\begin{lemma}[Finite change inequalities]\label{lem:reducible-change}
Suppose self-orthogonal local conditions on a fixed self-dual module are
changed at finitely many places.  Let \(b_0\) be the sum of the old local
dimensions modulo the intersections with the new conditions, and let
\(j_0\) be the rank of old Selmer localization in that sum.  Then
\[
 \dim\operatorname{Sel}_{\rm new}
 \le\dim\operatorname{Sel}_{\rm old}+b_0-2j_0.
\]
For a pair of dual scalar problems, suppose transverse lines are switched
at \(n\) odd places of local dimension two.  If their old localization
ranks are \(j,j_*\), then the first Selmer dimension increases by at most
\(n-j-j_*\), and the difference of the two Selmer dimensions is unchanged.
These statements include unramified old conditions at newly allowed
primes.
\end{lemma}

\begin{proof}
The quotients of the old and new conditions by their intersections pair
perfectly.  Their global localization images annihilate one another by
reciprocity.  The common strict kernel has dimension
\(\dim\operatorname{Sel}_{\rm old}-j_0\), while the new localization image
has dimension at most \(b_0-j_0\).  This proves the first inequality.
The same calculation for orthogonal scalar problems gives
\(n-j-j_*\).

For the difference assertion, compare each configuration with the one in
which the first local conditions are enlarged to the full local spaces.
Finite Poitou--Tate duality identifies the image of the enlarged global
Selmer group in the local quotient with the annihilator of the old dual
image.  The difference of the two Selmer dimensions therefore increases
by the sum of the local dimension increases.  That sum is the same for
either of the transverse-line configurations, proving the assertion.
These are the strict/relaxed localization calculations of finite duality;
see also \cite[Lemma~3.2 and Proposition~3.3]{MazurRubin10} and
\cite{MilneADT}.
\end{proof}

\subsection{The full rational plane}

Suppose first that \(W=E[2]\) is constant.  Prescribe \(\sum_i l_i\) for the
initial \(S\), and consider nonempty products \(h=\prod_i p_i^*\) with
that sum.  We choose the sum so that \(E^h\) has even functional sign.
The coprime twist sign formula permits this, for example by requiring
local squareclasses to be trivial at level primes and choosing the sign
of \(h\) appropriately.  These fixed local squareclasses make the Kummer
condition at \(S\) independent of the chosen primes.  By the parity
statement of Section~\ref{sec:classical}, every twist in this family has
even pure dimension.

Use additive notation for characters, with \(\mu_2\) identified with
\(\mathbb F_2\).  Kummer evaluations on halves of \(t\in W\) for the
untwisted curve are homomorphisms into \(W\), linear in \(t\) and
unramified outside \(S\).  Their value at a label \(l\) defines an operator
\(B(l)\); thus \(B\in H\otimes\operatorname{End}(W)\).

\begin{lemma}[The finite graph model]\label{lem:reducible-graph}
For \(h=\prod_{i=1}^n p_i^*\) in this family, every global two-Selmer class
has unique coordinates \(w\in H\otimes W\) and \(u_i\in W\).  They satisfy
\[
\begin{split}
 {\rm loc}_S(w)+\textstyle\sum_i \xi(l_i)\otimes u_i
 &\in L:=\bigoplus_{v\in S}\operatorname{im}(\delta_{2,E^h,v})
                     \ \subset V\otimes W,\\
 w(l_i)&=B(l_i)u_i+\sum_{j\ne i}e_{ij}(u_i+u_j)
                      \qquad(1\le i\le n).
\end{split}\tag{graph}\label{eq:graph}
\]
Here \(L\) is fixed throughout the family.  The torsion subspace is
\(u_i=t,\ w=Bt\), for \(t\in W\), and the dimension of the solution space
modulo that subspace is the pure dimension.
\end{lemma}

\begin{proof}
The radicals supported at \(S\) and the new signed primes give independent
coordinates on the global classes unramified outside their union.  This
gives the coordinates in the statement and its boundary condition at
\(S\).

At a new prime \(p_i\), the untwisted two-power torsion is unramified
because the curve has good reduction.  Twisting makes a ramified inertia
generator act by negation.  Its invariant two-power torsion is therefore
exactly \(W\).  The two-adic completion of the local points is this
torsion group: the remaining open subgroup is pro-\(p_i\).  On a half of
\(t\in W\), the twisted Kummer cocycle has inertia value \(t\) and
Frobenius value
\[
 B(l_i)t+\chi_h(\operatorname{Fr}_{p_i})t.
\]
Comparing with the global coordinates sets \(t=u_i\).  The contribution
of \(p_i^*\) itself occurs on both sides and cancels; the other radical
contributions give exactly the site equation in \eqref{eq:graph}.
The nonempty ramification set likewise forces the global two-primary
torsion to be exactly \(W\).  Its Kummer coordinates are
\(u_i=t,w=Bt\), as asserted.
\end{proof}

For a line \(A=\mathbb F_2a\subset W\), put \(\phi:E\to E/A\) and define
\(f_A\in H\) by
\[
 f_A(l)=B(l)a\bmod A,
\]
identifying \(W/A\) with \(\mathbb F_2\).  Also put
\(k_A=\dim(L\cap(V\otimes A))\).

\begin{lemma}[Orientation by isogeny]\label{lem:reducible-orientation}
We may replace \(E\) by a rationally isogenous curve with full rational
two-torsion, and apply that isogeny to all its twists, so that
\[
                    k_A\le m\qquad\text{for every }f_A=0.
                    \tag{orient}\label{eq:orient}
\]
The conductor and the prescribed family of local squareclasses are
preserved.
\end{lemma}

\begin{proof}
The image in \(E/A\) of a half of \(a\) is an additional two-torsion point.
It is rational exactly when \(B(-)a\) takes values in \(A\).  Thus
\(f_A=0\) exactly when \(E/A\) has full rational two-torsion.

The map \(H^1(A)\to H^1(W)\) is injective.  Projecting a two-Kummer class
to \(H^1(W/A)\) tests divisibility of the point by the dual isogeny;
inside the kernel, the classes are precisely the images of the
\(\phi\)-Kummer condition.  Consequently \(L\cap(V\otimes A)\) is the
sum of the local isogeny Kummer spaces.  Since the rational kernel has
order two at each place, \(k_A-m\) is the two-valuation of the product
over \(S\) of local cokernel orders divided by local kernel orders for
\(\phi:E^h\to(E/A)^h\).

When \(f_A=0\), this product equals the target/source ratio of
\(\Omega\prod_q c_q\).  Indeed, the product over all places has that
value by Haar change of variables: the finite N\'eron volumes are
\(c_q/L_q(E^h,1)\), the Euler factors are isogeny-invariant, and the
pullback scalar on differentials cancels by the product formula.  At
unused good odd primes the local ratio is one, by the unramified isogeny
Kummer condition.  At used primes \(p_i\), both two-primary point
completions have order four and the map has kernel of order two, so the
ratio is again one.  The product at \(S\) is fixed by the prescribed
local squareclasses.  This argument uses the total real period.

Consider the connected graph of full-rational curves joined to the
starting curve by full-to-full two-isogenies.  Every vertex is good
outside \(S\); Shafarevich finiteness makes this graph finite up to
rational isomorphism.  After one allowed twist, choose a vertex maximizing
the relative two-valuation of \(\Omega\prod_q c_q\); its ratios to the
other vertices are rational by the preceding calculation.  At every
outgoing full-to-full isogeny the target/source valuation is nonpositive,
which is \eqref{eq:orient}.  Isogenies preserve the conductor, and the
same inequality holds throughout the prescribed family.
\end{proof}

\begin{lemma}[Full-plane minimization]\label{lem:reducible-full-seed}
For the oriented curve of Lemma~\ref{lem:reducible-orientation}, some
allowed nonempty twist has pure dimension zero.  In particular,
Proposition~\ref{prop:reducible-seed} holds when the initial curve has full
rational two-torsion.
\end{lemma}

\begin{proof}
Choose a label list containing at least \(2m+2\) copies of every element
of \(H^*\).  Adjust multiplicities to obtain the prescribed total label.
Among all admissible edge assignments, minimize the pure dimension
\(d\) in \eqref{eq:graph}.  Lemma~\ref{lem:reducible-labels} realizes each
assignment by primes, so the minimum is even.  Suppose for a contradiction
that \(d\ge2\), and fix a minimizing graph.

We will vary only the edges.  Minimality first forces all differences
\(u_i+u_j\) of pure solutions into one line \(A\subset W\).  Edge
switches preserving that line then give more independent site errors
than the boundary condition at \(S\) permits.  The orientation inequality
\eqref{eq:orient} closes the argument when \(f_A=0\).

\paragraph{Even edge switches.}
For a nonzero even-weight vector \(v\in\mathbb F_2^n\), relax the site
equations by allowing errors \(v_i y\), with \(y\in W\).  The resulting
solutions map to \(W\oplus W\) by
\[
 (w,(u_i),y)\longmapsto(x,y),\qquad x=\sum_i v_i u_i.
\]
Let \(U\) be their image.  It is isotropic for the nondegenerate form
\[
 ((x,y),(x',y'))\longmapsto e(x,y')+e(y,x'),
\]
where \(e\) is the Weil form on \(W\).  To see this, apply global
reciprocity to two relaxed solutions.  At a changed site the error adds
an unramified class with Frobenius value \(v_i y\), while inertia remains
\(u_i\).  The cross terms are exactly the displayed form by local
Hilbert duality.  Kummer--Kummer and unramified--unramified terms vanish,
and every other local condition remains self-orthogonal.  Hence
\(\dim U\le2\).

The original graph imposes \(y=0\).  Toggle both directed edges on each
pair inside the support of \(v\).  Its even weight makes the change in
the \(i\)-th equation equal to \(v_i x\), so the toggled graph imposes
\(y=x\).  The two graphs share the strict kernel \(x=y=0\), which
contains their common torsion coordinates.  If \(x\) had rank two on the
old pure solution space, then \(U=W\oplus0\), and the toggled graph would
lose both those dimensions.  This contradicts minimality.  Thus every
even-sum map \(x\) has rank at most one on the pure solution space.

\paragraph{A common image line.}
These maps form a linear family and are jointly injective.  Indeed, if
all even sums vanish, the \(u_i\) are constant; subtracting the associated
torsion coordinates makes them all zero.  The site equations then give
\(w(l_i)=0\) at every label, whence \(w=0\).

A linear family of rank-at-most-one maps either has a common image line
or has a common domain functional: two rank-one tensors with both factors
independent have a rank-two sum.  The second possibility contradicts joint
injectivity on a space of dimension \(d\ge2\).  There is therefore a
common image line \(A=\mathbb F_2a\).  Pairwise even sums show that, after
subtracting torsion, every solution can be represented by \(u_i=t_i a\).
Projection of the site equations to \(W/A\) gives
\[
 \bar w(l_i)=t_i f_A(l_i).
\]
All labels occur, so \(\bar w\) vanishes on \(\ker f_A\) and is a multiple
of \(f_A\).  Subtracting the torsion coordinates in direction \(a\), when
necessary, gives representatives satisfying
\[
 w\in H\otimes A,\qquad t_i=0\quad\text{if }f_A(l_i)=1.
\]

\paragraph{The image line persists under the switch.}
Take an even \(v\) whose evaluation on the pure group is nonzero.  Its
strict pure kernel has dimension \(d-1\).  Since \(\dim U\le2\), and
\(U\) already contains a nonzero horizontal line, its intersection with
the diagonal has dimension at most one.  The toggled pure dimension is
therefore at most \(d\); minimality makes it exactly \(d\).  It is another
minimizing graph.  The shared strict pure space has dimension
\(d-1\ge1\).  Any nonzero vector in that space has some nonzero even-sum
evaluation, by joint injectivity, and that evaluation belongs to the image
lines of both graphs.  Thus their image line is the same \(A\).

The toggled graph consequently supplies a solution with nonzero
\(y=x=a\).  Normalizing it by torsion as above gives the same restrictions
on \(w,t_i\), while its error in the old graph is \((v_i a)_i\).

\paragraph{The dimension contradiction.}
Let \(D_A\) be the space of coordinates with these restrictions and the
boundary condition at \(S\), leaving the site errors unrestricted.  The
errors automatically lie in \(A\): their projection is
\(t_i f_A(l_i)=0\).  The error map from \(D_A\) for the original graph has
rank at least \(n-1\).  Indeed it contains \((v_i a)_i\) for every
nonzero-evaluating even \(v\), and these vectors span the even-weight
hyperplane.  To justify the latter assertion, the vanishing evaluations
form a proper linear subspace of that hyperplane, whose complement spans
the whole hyperplane.

The kernel of the error map has dimension \(d+i_T\), where
\(i_T=1\) if \(f_A=0\) and \(i_T=0\) otherwise.  These are precisely the
restricted torsion coordinates remaining after the pure quotient.  Put
\[
 n_1=\#\{i:f_A(l_i)=1\},\qquad r=m-\mathbf1_{f_A\ne0}.
\]
The allowable label sums \(\sum_i t_i l_i\) have rank \(r\), since the
allowable labels span \(\ker f_A\).  By
Lemma~\ref{lem:reducible-labels}, the boundary value is injective on the
pair \((w,\sum_i t_i l_i)\).  The kernel contributed by the \(t_i\) thus
has dimension \(n-n_1-r\), while the boundary image has dimension at most
\(k_A\).  We obtain
\[
 n-1+d+i_T\le\dim D_A\le n-n_1-r+k_A.
\]
If \(f_A\ne0\), the repeated labels and \(k_A\le\dim V=2m\) make this
impossible.  If \(f_A=0\), it gives
\(d\le k_A-m\le0\) by \eqref{eq:orient}, again a contradiction.
The minimum is therefore zero, proving the lemma.
\end{proof}

\subsection{The one-line seed}

\begin{lemma}[One-line construction]\label{lem:reducible-one-line}
Proposition~\ref{prop:reducible-seed} also holds when
\(\dim E(\mathbb Q)[2]=1\).
\end{lemma}

\begin{proof}
If a full-rational curve is available by a rational two-isogeny, use
Lemma~\ref{lem:reducible-full-seed} on that curve.  Otherwise first take
an even-functional-sign twist away from \(2N_E\), possibly the trivial
twist, using the same local prescriptions.  Denote the current curve by
\(E\) and enlarge \(S\) to include all current and original exclusions.
Choose a model
\[
 E:y^2=x(x^2+ax+b),\qquad D=a^2-4b,
\]
with kernel \(A=\langle T\rangle\), \(T=(0,0)\), and two-isogeny
\(\phi:E\to E/A\).  The standard quotient equation is
\[
 E/A:y^2=x(x^2-2ax+D).
\]
It follows from this equation, or the duplication law, that the remaining
two-torsion fields of the source and target are respectively
\(\mathbb Q(\sqrt D)\) and \(\mathbb Q(\sqrt b)\).  Thus \(D,b\) are both
nonsquare.  Write \(W=E[2]\) and \(\widehat\phi\) for the dual isogeny.
The module \(W\) is an extension of the scalar module \(W/A\) by \(A\)
with extension character \(D\).  Quadratic twisting preserves both
two-torsion modules and these squareclasses: a twist by \(h'\) replaces
\((a,b)\) by \((h'a,(h')^2b)\).

\paragraph{The isogeny Selmer groups.}
Let \(u,u_*\) be the dimensions of
\(\operatorname{Sel}_\phi,\operatorname{Sel}_{\widehat\phi}\), viewed as
radical groups supported at \(S\).  They contain the nontrivial classes
\(D,b\), respectively: these are the connecting classes of the opposite
kernel points, whose preimages are two-torsion.  For
\(F=\operatorname{Sel}_2(E)\), projection to \(W/A\) gives
\[
 \beta(F)\subset\operatorname{Sel}_{\widehat\phi},\qquad
 b\in\beta(F),\qquad
 F_A:=\ker(\beta|_F)\simeq\operatorname{Sel}_\phi/\langle D\rangle.
\]
These assertions follow from the Kummer diagrams for the factorization
of multiplication by two; \(\delta_2(T)\) gives the inclusion of \(b\).
For completeness, a lift of a kernel class to \(H^1(\mathbb Q,A)\) is
locally in the \(\phi\)-condition: its local two-Kummer point is divisible
by \(\widehat\phi\), and the ambiguity in the \(\phi\)-connecting class
is the image of the dual kernel under the same connecting map.  Globally
this ambiguity is \(\langle D\rangle\).  Even-sign parity and the one
rational two-torsion point imply that \(\dim F\) is odd at this stage and
at every subsequent stage.

\paragraph{The new local conditions.}
Each step twists by fresh signed primes whose labels sum to zero.
Lemma~\ref{lem:reducible-labels} makes the twisting product a local square
at all old places of \(S\), preserving the even sign.  Use only labels
with \(D(l_i)=0\).  The ramified-twist two-primary local point completion
is then the Frobenius-fixed two-torsion of each curve, by goodness and the
inertia-negation argument in Lemma~\ref{lem:reducible-graph}.

If \(b(l_i)=1\), the new \(\phi\)-condition is zero: the source completion
has order four, the target completion order two, and the map has kernel
of order two.  If \(b(l_i)=0\), the new conditions for both isogenies are
lines transverse to their old unramified lines.  In fact a nonkernel
two-torsion point in the target has a preimage in four-torsion under the
twisted isogeny with nonzero double.  Inertia acts by negation on this
preimage, producing a ramified Kummer class; the same argument applies
on the opposite side.  In either case the two-Kummer plane for \(W\) is
transverse to its old unramified plane, by the inertia calculation in
\eqref{eq:graph}.

\paragraph{Reducing the scalar Selmer dimensions.}
There are two cases.
\begin{enumerate}[label=\textup{(\roman*)}]
\item If \(D,b\) are independent squareclasses, use pairs of equal labels
with \(D(l)=0,b(l)=1\).  The new zero \(\phi\)-conditions impose vanishing
inside the old Selmer group, without admitting new ramified classes.
These labels span \(D^\perp\) on the current support.  Whenever \(u>1\),
one therefore evaluates nontrivially on
\(\operatorname{Sel}_\phi/\langle D\rangle\), and the corresponding pair
decreases \(u\).  This reaches \(u=1\).
\item If \(D=b\) as squareclasses, exchange the isogeny and its dual if
necessary so that \(u\le u_*\).  Equality of the squareclasses is
preserved.  While \(u\ge3\), choose labels
\(l,l',l+l'\in D^\perp\) with evaluation rank two on both isogeny Selmer
groups.  Such a choice exists: choose a two-plane modulo
\(\langle D\rangle\) in each group, prescribe isomorphisms from these
planes to \(\mathbb F_2^2\) agreeing on their intersection, and extend to
their sum and then to the full radical space.  The resulting two
functionals are \(l,l'\).

Both sets of local lines switch transversely.  By
Lemma~\ref{lem:reducible-change}, the difference \(u-u_*\) remains fixed
and
\[
 u_{\rm new}\le u+3-2-2<u.
\]
Thus the orientation persists, and iteration reaches \(u\le2\).
\end{enumerate}
Enlarge \(S\) after each step to contain the new primes.

\paragraph{A rank-three evaluation on the full Selmer group.}
Now suppose \(\dim F\ge3\).  Again use a repeated label pair \(l,l\).
In case (i), take \(D(l)=0,b(l)=1\), so \(u=1\) is maintained.  In case
(ii), the equality \(\dim F_A=u-1\) gives \(\dim\beta(F)\ge2\).  Choose
\(l\in D^\perp\) nonzero on \(\beta(F)\), and also nonzero on some
\(x\in\operatorname{Sel}_\phi\setminus\langle D\rangle\) if \(u=2\).
These at most two nonvanishing requirements modulo \(\langle D\rangle\)
are compatible over \(\mathbb F_2\).  At the two new places, the old dual
evaluation is nonzero, and the old primal evaluation is nonzero when
\(u=2\).  The transverse-line inequality therefore keeps \(u_{\rm new}\le2\).

We claim that Frobenius representatives for the pair can be chosen so
that evaluation of \(F\) in \(W\oplus W\) has rank at least three.
Their actions on \(W\) are trivial.  Their common projection to \(W/A\)
is the same nonzero functional \(\beta(-)(l)\); in case (i), its
nonvanishing follows from \(b\in\beta(F)\).  On the kernel of this
functional, which has dimension at least two, evaluations take values
in \(A\).  The subspace \(F_A\) has dimension at most one; if it has
dimension one, its diagonal evaluation in \(A\oplus A\) is nonzero by
the condition imposed using \(x\).

Here is the remaining freedom in those \(A\)-coordinates.  Let
\(\mathcal N\) be the kernel of all characters in \(H\), and restrict the
cocycles to \(\mathcal N\).  They become additive
\(A\)-valued characters there, and every class in \(F_A\) restricts to
zero.  Conversely, a class whose restriction vanishes factors through
the elementary quotient \(H^*\).  The square relation for its cocycle is
\[
 D(z)\,\beta(z)=0\qquad(z\in H^*),
\]
where \(\beta\) here denotes its projected character.  As \(D\ne0\),
this forces \(\beta=0\): two nonzero linear functionals cannot have
identically zero product.  The kernel of restriction is therefore exactly
\(F_A\).

It follows that the restrictions of \(F/F_A\) are independent additive
characters and realize arbitrary joint values.  Multiplying a Frobenius
representative of label \(l\) by elements of \(\mathcal N\) varies the
\(A\)-coordinate by any functional modulo \(F_A\), independently at the
two primes.  All data are finite and continuous, so Chebotarev realizes
these prescriptions at fresh primes outside \(S\).  If \(F_A=0\), choose
two independent functionals on the kernel of the
projected functional.  If \(\dim F_A=1\), choose two distinct extensions
of its fixed nonzero functional; they are independent over
\(\mathbb F_2\).  Thus the \(A\oplus A\)-evaluation has rank two on
that kernel.  Adding the common projected functional gives the claimed
total rank of at least three.

The self-dual change inequality at these two places now gives
\[
 \dim F_{\rm new}\le\dim F+4-2\cdot3=\dim F-2.
\]
The scalar Selmer bounds already achieved are retained.  Iterating reaches
\(\dim F=1\), since it is always odd and contains the torsion contribution.
This proves \eqref{eq:seed}, with the accumulated squarefree twisting
product and any chosen rational isogeny.
\end{proof}

\begin{proof}[Proof of Proposition~\ref{prop:reducible-seed}]
Apply Lemma~\ref{lem:reducible-full-seed} or
Lemma~\ref{lem:reducible-one-line}, according to the rational two-torsion
of the initial curve.  All twisting primes are fresh and avoid the
original \(2N_E\), so their product has the required discriminant and
coprimality properties.
\end{proof}

\subsection{A split imaginary partner and the analytic conclusion}

\begin{proposition}[Split imaginary Selmer partner]\label{prop:reducible-partner}
For \(h,E_0\) in Proposition~\ref{prop:reducible-seed}, there is a fresh
prime \(d\equiv7\pmod{16}\) such that
\(K=\mathbb Q(\sqrt{-d})\) splits all primes of \(2hN_E\), and
\[
 d_2(E_0^{-d})=s_2(E_0^{-d})=1,\qquad
 s_2(E_0)=0,\qquad s_2(E_0/K)=1.
\]
\end{proposition}

\begin{proof}
Resume notation \(E\) for the original curve and put \(W=E_0[2]\).
Choose \(S\) for \(E_0\) containing \(2,\infty\) and all primes of
\(hN_E\); in the one-line case also include the support of its final
model, whose coefficients we again denote by \(a,b,D\).
Choose a new prime \(d\equiv7\pmod{16}\) whose Frobenius acts as complex
conjugation on \(\mathbb Q(\sqrt H)\).  These prescriptions are compatible:
the elementary subfield of \(\mathbb Q(\mu_{16})\) is contained in
\(\mathbb Q(\mu_8)\), on which the exponent \(7\) acts as complex
conjugation.  We will add one further compatible condition in a case
below.  Signed-prime reciprocity makes \(-d\) a square at every odd finite
place of \(S\), and \(-d\equiv1\pmod8\) gives the same conclusion at \(2\).
The two-Kummer conditions of \(E_0\) and \(E_0^{-d}\) thus change
only at \(\infty,d\), and Frobenius at \(d\) acts on \(W\) as conjugation.
Quadratic twisting preserves the global rational two-torsion dimension,
so every change of Selmer dimension gives the same change of pure dimension.

\paragraph{Nontrivial conjugation on \(W\).}
If conjugation is nontrivial, it is the nontrivial Jordan block in
dimension two over \(\mathbb F_2\).  Hence
\(H^1(\mathbb R,W)=0\), while the old unramified condition at \(d\) has
dimension one.  Lemma~\ref{lem:reducible-change} bounds the increase of
the pure dimension by one.

\paragraph{Trivial conjugation on \(W\).}
Otherwise \(W\) is full locally at both places.  The planes at \(d\)
switch transversely.  At infinity the conditions are complementary lines
in \(H^1(\mathbb R,W)\simeq W\).  Indeed, on the two-adic Tate lattice,
complex conjugation \(c\) is identity modulo two and has determinant
\(-1\).  The rank-one projectors \((1+c)/2\) and \((1-c)/2\) are therefore
integral and give a direct decomposition.  Kummer evaluation on halves
of \(W\), namely \((c-1)/2\), spans the negative-eigenlattice line modulo
two; the negative twist spans the other line.  Local duality shows that
these are the full real Kummer conditions.

Thus \(b_0=3\) in Lemma~\ref{lem:reducible-change}.  Starting from
\eqref{eq:seed}, it suffices to obtain \(j_0\ge1\), or equivalently a
nonzero localization of an old global torsion two-Kummer class.  This
already occurs at infinity if \(W\) is globally full.  It also occurs if
the unique rational kernel point \(T=(0,0)\) lies outside the real
identity component.

\paragraph{The remaining quartic prescription.}
In the remaining one-line case, \(T\) is the greatest root of the real
cubic, whose component contains infinity.  Hence \(a,b,D>0\) and
\(a>2\sqrt b\), with
\(b,D\) nonsquare.  Prescribe in addition that Frobenius at \(d\) flip
\(\sqrt{a+2\sqrt b}\), while fixing \(\sqrt b,\sqrt D\), as the earlier
prescription already requires.  The duplication law gives the halves
of \(T\) as
\[
 x=\pm\sqrt b,\qquad y^2=b(a\pm2\sqrt b).
\]
Since \(\sqrt D\) is fixed, the two conjugate square roots flip
simultaneously.  The corresponding half \(Q\) is sent to \(-Q\), so
its Kummer value is \(-2Q=T\ne0\).  Thus the additional prescription
gives the required nonzero torsion localization at \(d\).

To verify compatibility, let \(L\) be the splitting field generated by
\(\sqrt b,\sqrt D,\sqrt{a+2\sqrt b}\).  It contains the conjugate square
root by the norm relation.  Put \(P=\mathbb Q(\sqrt H,\mu_{16})\).
If \(b,D\) are independent squareclasses,
the two conjugate radicands over \(\mathbb Q(\sqrt b)\) are independent
squareclasses: their norm is \(D\), and their product \(D\) is still
nonsquare in that field.  Conjugacy exchanges the independent sign
actions, so their simultaneous flip is a commutator.  It therefore fixes
\(L\cap P\) and can be adjusted freely over
the prescribed abelian data, which already fix \(\sqrt b,\sqrt D\).

If \(D=b\) as squareclasses, the radicand is still nonsquare by its norm,
and \(L\) is cyclic quartic.  A lift of the nontrivial automorphism of
\(\mathbb Q(\sqrt b)\) also flips \(\sqrt D\), so its square is the
simultaneous flip.  The unique quadratic subfield of \(L\) is
\(\mathbb Q(\sqrt b)\).
If \(L\not\subset P\), then \(L\cap P=\mathbb Q(\sqrt b)\), so the
central flip can be chosen freely over the already prescribed data.
In the containment case, \(L\) is totally real because
\(a\pm2\sqrt b>0\).  The square subgroup of
\(\operatorname{Gal}(P/\mathbb Q)\) has order two: the radical extension
has exponent two, and the cyclotomic factor contributes the square
acting by exponent \(9\) on \(\mu_{16}\).  The prescribed Frobenius on
\(P\) differs from complex conjugation by this nontrivial square.
Every cyclic quartic quotient has a nontrivial square subgroup, so this
element acts nontrivially on \(L\).  Since conjugation is trivial on the
totally real field \(L\), the prescribed Frobenius gives exactly the
required flip.  Chebotarev now supplies \(d\) satisfying all prescriptions.

We have proved in every case that \(d_2(E_0^{-d})\le1\).  The coprime
split imaginary twist changes the functional sign, so pure-dimension
parity makes this dimension exactly one.  The finite-Selmer bound and
corank parity then give \(s_2(E_0^{-d})=1\), while \eqref{eq:seed} gives
\(s_2(E_0)=0\).  Finally, quadratic decomposition of the ordinary Selmer
coranks gives \(s_2(E_0/K)=1\).
\end{proof}

\begin{corollary}[The reducible split-pair anchor]\label{cor:reducible-anchor}
Proposition~\ref{prop:anchor} holds for every non-CM elliptic curve over
\(\mathbb Q\) with rational two-torsion.
\end{corollary}

\begin{proof}
Apply the reducible split-product comparison of
Sections~\ref{sec:split-measures} and~\ref{sec:split-local} to \(E_0,K\)
from Proposition~\ref{prop:reducible-partner}.  The curve remains non-CM,
its level and \(2\) split in \(K\), and its ordinary Selmer corank over
\(K\) is one.  That comparison gives a simple zero of the product and,
by \eqref{eq:E} and \eqref{eq:G},
\[
 \operatorname{an}(E_0)+\operatorname{an}(E_0^{-d})=1,
 \qquad X(E_0)+X(E_0^{-d})=0.
\]
At this point both factors have analytic rank at most one, so the
classical theorem gives finiteness of their Tate--Shafarevich groups.
We may therefore use the isogeny arithmetic-volume comparison of
Section~\ref{sec:classical} to transfer these conclusions from \(E_0\)
and \(E_0^{-d}\) to \(E^h\) and \(E^{-hd}\).

Every signed-prime product used to form \(h\) was squarefree, odd, and
coprime to \(2N_E\), with the empty product allowed.  The fresh prime
\(d\) avoids its support, and \(K\) splits all primes of \(hN_E\) as
well as \(2\).  Thus \(k=-d\) and \(h\) satisfy all discriminant,
coprimality, and splitting requirements of Proposition~\ref{prop:anchor}.
\end{proof}
